\section{Diagnosis: a map that detects, ranks, and spreads}
\label{sec:diagnosis}
\begin{table}[t]
\centering
\caption{\textbf{What the released SALAD map gives an inspector.} Extracted path, 138 test trays. $R$: target recall. $W$: intact compartments queued per faulty tray (13 per relocation, 14 per vacancy). Pooled trays mix relocations (two targets) and vacancies (one), with every parameter refit on pooled development trays. Clean trays: compartments queued per unedited source and per valid sham. The rule uses the published constraint and is extra task knowledge.}
\label{tab:released}
\footnotesize\setlength{\tabcolsep}{2.4pt}
\begin{tabular}{lrrrrc}
\toprule
Queue policy & \multicolumn{2}{c}{Relocations} & \multicolumn{2}{c}{Pooled} & Clean trays \\
\cmidrule(lr){2-3}\cmidrule(lr){4-5}
 & $R$ (\%) & $W$ & $R$ (\%) & $W$ & Src.\,/\,sham \\
\midrule
Threshold, 95\% dev.\ recall & 95.7 & 11.07 & 93.7 & 4.58 & 0.00 / 0.00 \\
Relative, 95\% dev.\ recall & 95.7 & 11.07 & 93.7 & 4.58 & 0.04 / 0.07 \\
List, $k$ fixed on dev. ($k{=}3$/3) & 97.1 & 1.06 & 98.1 & 1.53 & 0.11 / 0.15 \\
Top 2 (fixed) & 88.8 & 0.22 & 92.5 & 0.61 & -- \\
Top 3 (fixed) & 97.1 & 1.06 & 98.1 & 1.53 & -- \\
\midrule
Explicit rule (extra knowledge) & 100.0 & 0.00 & 100.0 & 0.00 & 0.00 / 0.00 \\
\bottomrule
\end{tabular}

\end{table}

\paragraph{Detection and ranking are excellent.}
The released composition branch alarms on all \relFaults{} relocations and all \relVacancies{} vacancies, and scores every fault above its valid sham. On the oracle path the two affected compartments rank at \relAUC\% AUROC after the edit, against \relAUCsrc\% for the same compartments in the unedited source, so the ranking is caused by the edit and not inherited from the scene. Only \relSrcAlerts{} unedited sources and \relShamAlerts{} valid shams alarm on the extracted path.

\paragraph{Thresholding is not.}
\Cref{tab:released} shows what this means for an inspector. The cutoff that reaches 95\% target recall on development trays opens \relThrW{} intact compartments per faulty tray on test, at \relThrR\% recall. Scaling the cutoff to each tray's maximum does not help, because faulty trays saturate near one. A list of the top three compartments, the shortest list reaching 95\% on development, opens \relListW{} intact compartments at \relListR\% recall. The released map therefore ranks well and calibrates poorly: what separates the targets from intact compartments is their order within a tray, not their absolute score. Lists exploit that order only when the number of faulty compartments is known. Pooling relocations with vacancies, the list must stay at three and wastes \relPoolListW{} compartments per tray. The published rule, applied to the same extracted part maps, finds every target with no waste; we treat it as an upper bound that requires knowing the rule.

\paragraph{The controls locate the spread.}
\Cref{fig:spread} compares each edited map with its source on the same 13 untouched compartments. After a relocation, \relAnyFault\% of them contain a pixel that moves by at least 0.5 (\relAnyFaultX\% on the extracted path). After the vacancy-only control, which empties the same compartment, the figure is \relAnyVac\%. After the valid sham, which moves the same pin the same distance, it is \relAnySham\%. The response therefore follows the \emph{inserted} pin, not the vacancy or the edit. It is peaked rather than diffuse: \relArea\% of the untouched area crosses 0.5, mostly on the pins themselves (\cref{fig:teaser}).

\begin{figure}[t]
\centering
\includegraphics[width=\columnwidth]{spread.pdf}
\caption{\textbf{The spread follows insertion and class.} Left: fraction of the 13 untouched compartments whose map changes by more than a cutoff, for the three matched edits (oracle path, 95\% bootstrap bands). Right: after one extra instance is pasted into a part map, fraction of untouched instances of the \emph{same} class versus pixels of \emph{other} part classes that change by at least 0.5, for three released checkpoints.}
\label{fig:spread}
\end{figure}

\paragraph{The spread follows the inserted class.}
If the spread follows insertion, an inserted part should raise the other parts of its class and nothing else. We test this on three released checkpoints by pasting one intact donor instance from a training map onto the substrate on which that class rests, in map space, and scoring the composition branch (\cref{fig:spread}, right; per-category numbers in the supplement). In pushpins, an extra pin raises \ppAny\% of the untouched pins and 0.1\% of other-class pixels. In screw bags, an extra screw, nut or washer raises \sbAny\% of the untouched parts and no other-class pixels. In breakfast boxes, a fruit placed in the cereal compartment raises the original fruit in \bbAny\% of trays. The signature is strong for the two classes with many small instances and absent for the merged fruit region.

\section{Cause: the spatial extent of the insertion label}
\label{sec:cause}
\begin{table*}[t]
\centering
\caption{\textbf{One protocol for every detector.} Matched trays, extracted path: fault/vacancy and valid-sham alarms (of 138), spread (untouched compartments changing by $\ge0.5$ after a relocation; probability maps only), wasted inspections $W$ per faulty tray with target recall $R$ for the threshold and the development-selected list, $W$ for the threshold on pooled relocation and vacancy trays, compartments queued per valid sham, and the cell analogue of low-FPR sPRO with only faulty-tray negatives (fault) or also the source and sham compartments (all). Native LOCO: AU-sPRO up to 5\% FPR and image AUROC on the 91 logical anomalies and 138 good pushpin images. Class-wide and footprint heads are retrained from the released head on identical inputs (\cref{sec:cause}).}
\label{tab:main}
\footnotesize\setlength{\tabcolsep}{3.1pt}
\begin{tabular}{llrrrrrrrrrr}
\toprule
\multicolumn{2}{l}{Detector, insertion label} & \multicolumn{2}{c}{Alarms} & Spread & \multicolumn{2}{c}{Threshold} & \multicolumn{2}{c}{List} & Pooled & \multicolumn{2}{c}{Region area} \\
\cmidrule(lr){3-4}\cmidrule(lr){6-7}\cmidrule(lr){8-9}\cmidrule(lr){11-12}
 & & fault & sham & (\%) & $W$ & $R$ & $W$ & $R$ & $W$ & fault & all \\
\midrule
SALAD & released & 138 & 7 & 89.4 & 11.07 & 95.7 & 1.06 & 97.1 & 4.58 & 74.0 & 81.6 \\
 & class-wide, seed 1 & 138 & 19 & 33.8 & 7.72 & 93.5 & \textbf{1.04} & 97.8 & 2.85 & 74.9 & 83.0 \\
 & class-wide, seed 2 & 138 & 15 & 66.5 & 10.87 & 95.3 & 4.04 & 97.8 & 5.34 & 56.8 & 62.8 \\
 & footprint, seed 1 & 138 & 25 & 1.6 & \textbf{4.57} & 96.0 & 2.06 & 97.1 & \textbf{1.09} & 76.2 & \textbf{87.3} \\
 & footprint, seed 2 & 138 & \textbf{6} & \textbf{1.4} & 6.15 & 96.4 & 3.07 & 96.7 & 3.36 & 69.9 & 81.8 \\
\midrule
EfficientAD & -- & 138 & 128 & -- & 0.21 & 97.1 & 0.04 & 97.8 & 0.22 & 92.9 & 69.4 \\
PatchCore & -- & 133 & 123 & -- & 0.82 & 92.8 & 1.96 & 94.9 & 12.50 & 84.3 & 66.4 \\
\midrule
Explicit rule & (reference) & 138 & 0 & -- & 0.00 & 100.0 & -- & -- & 0.00 & -- & -- \\
\bottomrule
\end{tabular}

\end{table*}

\paragraph{The released label.}
SALAD trains its discriminative head on four synthetic edits of the part map: Perlin-shaped relabeling, relabeling of a connected component, insertion of a rotated component from another training map, and relabeling plus insertion \cite{fucka2025}. For insertions, the pinned implementation pastes the donor component and then returns the \emph{entire} channel of its class as the positive target (\texttt{copy\_label\_diff\_img}). When the pasted component is a pin, all 16 pins become positive. In 204 insertions replayed from the pinned generator on pushpin maps, the median target covers 19{,}800 pixels, 16 times the 416 pixels actually pasted. SALAD motivates this choice with interchangeable screws \cite{fucka2025}. It predicts exactly the class-specific spread of \cref{sec:diagnosis}.

\paragraph{A controlled retraining.}
We continue training the released discriminative U-Net on SALAD's own edit stream, keeping the released reconstruction autoencoder frozen: 1{,}600 steps of one possibly edited map plus its clean copy, the four pinned edit families with probability 1/2, five times sigmoid focal plus L1 loss, and Adam at $10^{-5}$ with a tenfold decay at 95\%. For each seed, one process copies the released weights into two arms and feeds both the same prefetched stream of maps, donors and edits, so the inputs are identical bit for bit. The arms differ only in the target of an insertion: the whole class channel (\emph{class-wide}, as released) or the pasted donor footprint (\emph{footprint}). Perlin and relabel targets are shared, and no rule-specific edit is used. Two seeds give two matched pairs. The supplement reports four further pairs from two earlier recipes, including one with rule-aware training edits.

\paragraph{Footprint labels remove the spread.}
\Cref{tab:main} compares the arms on the matched trays. Every head still alarms on all 138 faults and vacancies. With class-wide labels, \cwSpread\% of untouched compartments change by at least 0.5 after a relocation; with footprint labels, \fpSpread\%, a paired reduction of \dSpread{} points in the two seeds. The vacancy and sham responses stay at or below 3.1\% in every head. The same holds in all four supplementary pairs. The label alone therefore decides whether an inserted part lights up its whole class.

\paragraph{Removing the spread is not yet a short queue.}
At the 95\%-recall threshold, footprint labels cut wasted inspections by \dThrW{} compartments relative to their class-wide pairs (paired intervals; the same sign holds in both test halves). The footprint heads still open \fpThrW{} intact compartments per tray, because a short continuation leaves the two targets less separated from the rest: both are in the top three on \fpTopThree\% of trays, against \relTopThree\% for the released map. On matched relocations alone, a top-three list on the released map remains the cheaper policy (\relListW). When relocations and vacancies are pooled, the footprint threshold opens \fpPoolW{} compartments per tray, against \relPoolListW{} for the released list. Adding rule-aware training edits to the footprint label shortens the threshold queue to 0.20 and 0.45 (supplement). The label controls \emph{where} the map responds; how sharply it ranks the targets also depends on training.

\section{Real logical anomalies}
\label{sec:native}
\begin{table*}[t]
\centering
\caption{\textbf{Native MVTec LOCO logical anomalies} (official test split). AU-sPRO follows the official rules (per-defect saturation; false positives pooled over good and anomalous images) at $256{\times}256$. Two seeds per label policy. Rows below ``all logical'' restrict the sPRO average to one defect type.}
\label{tab:native}
\small\setlength{\tabcolsep}{5pt}
\begin{tabular}{llrrrrrcc}
\toprule
\multicolumn{2}{l}{Native test set} & Released & \multicolumn{2}{c}{Class-wide} & \multicolumn{2}{c}{Footprint} & \multicolumn{2}{c}{Footprint $-$ class-wide} \\
\cmidrule(lr){4-5}\cmidrule(lr){6-7}\cmidrule(lr){8-9}
 & & & seed 1 & seed 2 & seed 1 & seed 2 & seed 1 & seed 2 \\
\midrule
Pushpins & all logical & 33.9 & 34.9 & 33.4 & \textbf{78.5} & 71.6 & +43.6 [26.6, 56.0] & +38.2 [16.5, 54.7] \\
 & extra pin & 16.4 & 14.0 & 18.7 & \textbf{83.7} & 70.1 & +69.7 [47.4, 79.7] & +51.4 [21.0, 69.1] \\
 & missing pin & 45.2 & 47.7 & 41.7 & \textbf{70.8} & 69.6 & +23.2 [5.7, 41.9] & +27.9 [9.4, 48.5] \\
 & missing separator & 72.3 & 85.6 & 77.0 & 94.0 & \textbf{95.3} & +8.4 [5.6, 11.4] & +18.3 [10.9, 25.1] \\
\midrule
Breakfast box & all logical & 18.3 & 18.2 & 18.5 & 27.2 & \textbf{30.7} & +9.0 [4.2, 16.2] & +12.1 [6.9, 20.2] \\
 & missing cereals & 55.1 & 60.0 & 54.7 & 86.3 & \textbf{89.2} & +26.4 [11.6, 43.3] & +34.5 [18.3, 52.7] \\
\midrule
Juice bottle & all logical & 33.6 & \textbf{39.5} & 38.3 & 39.4 & 38.4 & $-$0.1 [$-$1.4, 1.9] & +0.0 [$-$1.5, 2.7] \\
\midrule
Splicing conn. & all logical & 17.8 & 18.0 & \textbf{23.6} & 18.8 & 18.5 & +0.8 [$-$0.7, 4.8] & $-$5.1 [$-$11.9, 1.0] \\
\midrule
Screw bag & all logical & \textbf{17.6} & 15.1 & 14.4 & 15.4 & 15.6 & +0.3 [$-$5.0, 2.6] & +1.2 [$-$2.1, 2.4] \\
 & extra part & 2.5 & \textbf{4.3} & 0.6 & 0.0 & 0.0 & $-$4.3 [$-$16.8, $-$0.1] & $-$0.6 [$-$9.6, $-$0.0] \\
\bottomrule
\end{tabular}

\end{table*}
The matched trays isolate the mechanism. The native test set shows whether it matters. MVTec LOCO annotates each logical defect region and scores localization by sPRO: each region counts as found once a defect-specific fraction of it is covered, and false positives are pixels outside every region, pooled over good and anomalous images \cite{bergmann2022}. We re-implemented the official rules, checked against the released evaluation code, on the $256{\times}256$ maps, and integrate up to 5\% false-positive rate.

\begin{figure}[t]
\centering
\includegraphics[width=\columnwidth]{native_examples.pdf}
\caption{\textbf{Native logical anomalies} (first test image of each type; cyan: ground truth). Each map is divided by that head's pixel threshold at 1\% false-positive rate on the 138 good test images, so bright pixels are those it would flag at that operating point. For the missing pin and the missing separator, the released and class-wide heads flag 64--76\% of all non-defect pixels and the footprint head 2--11\%, while each covers at least 89\% of the defect. For the extra pin, the released head flags the extra pins together with every other pin; at this strict threshold the footprint head flags little of either.}
\label{fig:native}
\end{figure}

On the 91 native logical pushpin anomalies, footprint labels raise AU-sPRO from \natPPcw\% for the class-wide pairs (\natPPrel\% released) to \natPPfp\% (\cref{tab:native}). The result does not depend on the recipe. Across all six matched pairs, from three training recipes (four of them in the supplement, including rule-aware training), footprint labels give \sixFP\% and class-wide labels \sixCW\%; the smallest paired gain is \sixMinGain{} points. The gain is largest where the mechanism predicts it. On extra-pin defects, where the released map lights up every pin, AU-sPRO rises from \natAddCW\% to \natAddFP\%. On missing pins it rises from \natMissRel\% (released) to \natMissFP\%. Here the class-wide maps flag most of the tray floor (\cref{fig:native}). This is consistent with the same mechanism: a vacated pin or separator becomes floor, and a class-wide label for an inserted floor region covers the entire floor. Image-level detection pays a small price: AUROC is \aucFP\% with footprint labels against \aucCW\% for the class-wide pairs. Both short continuations sit below the released head (\aucRel\%), which was trained jointly for 70{,}000 steps. These are real faults photographed by the dataset authors; no edit of ours is involved. \Cref{fig:native} shows the first test image of each defect type at a low-FPR operating point.

\paragraph{When the location is not identifiable.} Screw bags are SALAD's own motivating case: the rule counts interchangeable screws, nuts and washers, so no single instance is ``the'' extra one. MVTec LOCO annotates it that way. For an extra nut, the ground truth covers every nut in the bag and saturates once one of them is found, whereas for an extra pushpin it covers only the pins in the overfilled compartment. Here footprint labels bring nothing. Over all screw-bag logical anomalies, AU-sPRO is \natSBfp\% with footprint labels, \natSBcw\% with class-wide labels and \natSBrel\% released, and on extra-part defects the class-wide heads reach up to \natSBaddCW\% against \natSBaddFP\% (\cref{tab:native}). The benchmark's own annotation thus says where a label should point. The spatial extent of a synthetic label should follow the inspection task's definition of where the fault is, not a single convention for all classes.

\section{What matched shams reveal}
\label{sec:system}
\begin{table}[t]
\centering
\caption{\textbf{Full SALAD on matched pushpin trays.} Alarms out of 138 on the extracted path. Native: logical AUROC on 91 native logical anomalies and 138 good trays. The validation cutoff is the published operating point. The normal-test cutoff is refit on one fixed half of the unedited test trays and applied to the other half, in both directions.}
\label{tab:system}
\footnotesize\setlength{\tabcolsep}{1.7pt}
\begin{tabular}{lrrrrrrr}
\toprule
Branch & Native & \multicolumn{3}{c}{Validation cutoff} & \multicolumn{3}{c}{Normal-test cutoff} \\
\cmidrule(lr){3-5}\cmidrule(lr){6-8}
 & AUROC & src. & sham & fault & src. & sham & fault \\
\midrule
Fused & 100.0 & 75 & 130 & 138 & 13 & 64 & 138 \\
Local app. & 98.5 & 9 & 128 & 138 & 11 & 124 & 138 \\
Composition & 99.7 & 5 & 7 & 138 & 11 & 12 & 138 \\
Gaussian & 84.4 & 84 & 102 & 106 & 8 & 18 & 20 \\
\bottomrule
\end{tabular}

\end{table}

\paragraph{Appearance detectors find the edit.}
EfficientAD (with SALAD's released appearance weights) and PatchCore read RGB rather than part maps. On faulty trays alone they localize the targets well. Their fault-only region areas are \eadRegFault\% and \pcRegFault\%, against \relRegFault\% for SALAD's composition map, and EfficientAD's threshold queue opens only \eadThrW{} intact compartments (\cref{tab:main}). The valid shams tell a different story. EfficientAD alarms on \eadSham{} and PatchCore on \pcSham{} of 138 valid shams, against \eadSrc{} and \pcSrc{} unedited sources, and both queue the sham's edited compartment. Counting source and sham compartments as negatives reverses the ranking: \eadRegAll\% and \pcRegAll\% against \relRegAll\%. These maps locate where pixels were composited, which on faulty trays coincides with the targets. A fault-only evaluation, including sPRO on edited images, cannot tell the two apart.

\paragraph{Full SALAD.}
Full SALAD adds a local appearance branch (EfficientAD \cite{batzner2024}) and a global Gaussian branch over pooled student features \cite{rippel2021}, and sums their normalized scores with the composition score \cite{fucka2025}. It reaches 100\% logical AUROC on the native pushpin anomalies. With its validation-fitted cutoff, however, it alarms on \fullSrcVal{} of 138 unedited test trays (\cref{tab:system}). The Gaussian branch is responsible: its median normalized score is near zero on validation but 2.7 on unedited test trays. Refitting the cutoff on held-out normal test trays restores a \fullSrcX/138 source rate, which separates this calibration shift from the edits. The valid shams then expose a second problem. The fused score still alarms on \fullShamX{} of 138 valid shams, because the local appearance branch alarms on 124 of them and on only 11 sources. That branch sees the composite rather than the arrangement. The composition branch, which reads only the part map, alarms on 12 shams and 11 sources. Without the sham, all 138 fault alarms would be credited to logical reasoning.

\paragraph{A second rule.}
In juice bottles the upper label icon must match the liquid \cite{bergmann2022}. On 89 held-out bottles we paste an intact icon from a fixed training bank into the upper label: the same fruit is a valid sham, another fruit is a fault, with identical support, offset and blending. The composition branch cannot see the icon and ties exactly. The global appearance branch scores all 89 faults above their shams and alarms on all faults and 14 shams; the local appearance map alarms on 25 faults and 28 shams. Here the matched sham shows which branch carries the logical evidence (supplement).

\section{Limitations and conclusion}
\label{sec:conclusion}
The matched edits are composites of real photographs. They cover one kind of logical fault, and the extracted path still carries compositing traces that the oracle path removes. The inspection protocol needs compartments with an exact action label, which among the LOCO categories only pushpins provide; the insertion probe extends the diagnosis to screw bags but not the inspection cost. Our retraining continues the released head with a frozen autoencoder; it is not SALAD's full joint training. We evaluate the composition branch, two seeds per recipe, and our AU-sPRO is a re-implementation at map resolution. Class-wide labels have a purpose when the faulty instance is truly ambiguous. Our point is that this choice decides where the map points, and it should be made, and evaluated, for the localization task at hand.

Detection AUROC and threshold-free curves describe what a map can rank. Matched rearrangements describe what it asks an inspector to do, and why. For a state-of-the-art logical detector, the answer came down to the spatial extent of one synthetic label. Labelling the footprint that was actually pasted, which every synthetic edit already knows, more than doubled localization of real logical anomalies.
